
\scanpage{50}
\mspage{6}
\uline{Remarque} En fait, on a un foncteur plus précis, vers une catégorie de graphes munis d'une structure où la partie $\rho$ de la pondération est \deleted{\unsure{complétée}} \inserted{remplacée} par une structure plus précise, à savoir \deleted{une \unsure{application}} \deleted{\inserted{une \unsure{famille}}} famille
\[
 \deleted{\unsure{\alpha}\,(\widetilde R,\alpha,p:R\longrightarrow\vec A)}
 \qquad (R_a)_{a\in A}
\]
\deleted{\diagram{\begin{tikzcd}[ampersand replacement=\&,row sep=small]\vec A\arrow[d,two heads]\\ \vec A/\sigma\end{tikzcd}}}
d'ensembles, et pour $\forall a\in A$, si nous désignons l'\unsure{ens} des
\deleted{\unsure{donné d'une permutation de}}
``orientations'' de $a$ ($\operatorname{card}\omega(a)\in\{1,2\}$), une application
\deleted{$R$, et $p$ une application $R\longrightarrow A$ \unsure{conj. à}}
\[
 \omega\longmapsto u_\omega:\quad \omega(a)\longrightarrow\operatorname{Circ}(R_a)
\]
ens des permutations circulaires sur $R_a$
\deleted{et \unsure{d'autre} \unsure{une} \unsure{application injective}}
\[
 \deleted{\unsure{R/\alpha}\lhook\joinrel\longrightarrow\vec A}
\]
avec la condition que si $\omega,\omega'\in\omega(a)$, $\omega\ne\omega'$, alors
\[
 \deleted{u_\omega u_{\omega'}=\operatorname{id}_R}
 \qquad u_{\omega'}=u_\omega^{-1}.
\]
\margin{$\omega(a)\mathrel{\mathop=^{\rm def}}\{\vec a\in\vec A\mid\text{au dessus de }a\}$.}

On peut aussi considérer que la donnée des \unsure{suppléments} revient à la donnée de
\[
 \begin{array}{rcl}
 r:A&\longrightarrow&\mathbb N\\
 a&\longmapsto&\operatorname{Card}p^{-1}(a)
 \end{array}
\]
(la ``pondération-arêtes''), et de donnée d'une famille de torseurs
\[
 (R_a)_{a\in A},\quad \deleted{\unsure{\omega(a)}}
\]
$R_a$ un torseur sous $\mathbb Z/r(a)\mathbb Z\mathbin{\wedge}_{\{\pm1\}}\omega(a)$.
On associe à $(X,\widetilde\Sigma,R)$ une telle structure supplémentaire en prenant
\[
 R_a=R\cap\widetilde\Sigma_a
\]
($\widetilde\Sigma_a$ la comp. connexe de $\widetilde\Sigma$ \deleted{\unsure{marquée}} \inserted{qui correspond à} $a\in A\simeq\pi_0(\widetilde\Sigma)$),
\deleted{$\unsure{R=R(\Sigma C)|R=\pi_0(\widetilde\Sigma)}$, \unsure{$R\longrightarrow S$} \unsure{la donnée de} $\widetilde\Sigma$}
\deleted{\diagram{\begin{tikzcd}[ampersand replacement=\&,row sep=small,column sep=large]
R\arrow[r]\arrow[d,equal]\&S\arrow[d,equal]\\
\unsure{\Gamma}\arrow[r]\&\pi_0(\widetilde\Sigma)
\end{tikzcd}}}
$\omega\longmapsto u_\omega$ défini de façon évidente.
\deleted{et $p$ \unsure{déduit de} \unsure{$(\widetilde\Sigma_a)$} \unsure{sur} $R_a$.}

\uline{NB} Même si $r(\unsure{\omega|a})=1$ ou $r(\unsure{\omega|a})=2$, le fait de disposer
\mspage{49}
